\chapter{Los lenguajes de este libro}
\label{cap:lenguajes}

Hay una pregunta que en este libro se plantea en cada línea: \emph{esto que estoy leyendo,
¿se afirma dentro del sistema del que hablamos, o desde fuera?} Contestarla mal es caro — la Parte~\ref{parte:noescrito} cuenta una inconsistencia que fue exactamente
eso, un error de categoría. Este capítulo la contesta de una vez, y con más detalle del que suele
darse: porque ni hay una metateoría, ni hay un solo lenguaje del que se hable. Hay tres y hay
cuatro.

\section{Tres metateorías, no una}

Cuando se dice «la metateoría» se está diciendo, en tres momentos distintos de esta historia, tres
cosas incompatibles.

\textbf{La de Gödel era deliberadamente débil.} Razonamiento finitario sobre secuencias de símbolos,
y nada más — porque el objetivo era convencer a Hilbert, y una demostración de consistencia que
usara el infinito no habría valido de nada. Esa autolimitación no era humildad ni convicción
—§\ref{sec:finitismo-estrategia} cuenta que ninguno de los dos era constructivista—: era el
requisito para que el argumento le sirviera a quien iba dirigido.

\textbf{La del kernel de este proyecto es más fuerte de lo que parece.} El cálculo $\vdash$ de
\modulo{FOL} incluye la \defterm{$\omega$-regla}: de «$A[n]$ para todo término $n$» concluye
$\forall A$. Es válida en el modelo estándar y tiene \emph{infinitas premisas}, luego el cálculo que
la admite no es \defterm{r.e.} — no hay programa que liste sus teoremas.

\begin{muro}[La consecuencia que gobierna el libro entero]
Con la $\omega$-regla, $\vdash$ no es recursivamente enumerable. Y por el teorema de Tarski, ninguna
fórmula puede representar fielmente la demostrabilidad de un cálculo así. \textbf{Por eso el
proyecto construye un segundo cálculo}, finitario, y aritmetiza ése. Todo lo que en este libro se
codifica, se verifica o se diagonaliza ocurre en el segundo, no en el primero.
\end{muro}

\textbf{Y la de Lean es enormemente más fuerte que ambas.} Su teoría de tipos interpreta a ZFC y
más. Ahora bien —y esto sí está medido sobre el árbol—, \textbf{el proyecto usa dos niveles de ella
y ninguno más}: proposiciones y datos. No hay universos por encima, no hay polimorfismo de
universos, y no se cuantifica sobre proposiciones arbitrarias en ninguna parte. La fuerza está
disponible; no se gasta.

\subsection*{Dónde cae cada una: la escala, con nombres}

«Más fuerte» y «más débil» dicho así no informa. Hay una escala estándar, y las tres piezas de este
libro tienen sitio en ella. De arriba abajo por lo que \emph{demuestran} de aritmética:

\begin{center}\small
\begin{tabularx}{\linewidth}{@{}l >{\raggedright\arraybackslash}X l@{}}
\toprule
sistema & qué permite & papel aquí \\
\midrule
PRA & recursión primitiva, inducción sin cuantificadores no acotados & el finitismo que pedía Hilbert \\
$\mathrm{I}\Sigma_1$ & inducción restringida a fórmulas $\Sigma_1$ & lo que \emph{basta} como metateoría (cap.~\ref{cap:objeto-y-meta}) \\
HA & PA con lógica intuicionista & equiconsistente con PA \\
PA & inducción sobre todas las fórmulas de primer orden & la referencia habitual \\
Q\texttt{++} $+\ \omega$-regla & el cálculo $\vdash$ de este proyecto & \textbf{toda la aritmética verdadera} \\
\bottomrule
\end{tabularx}
\end{center}

La última fila no es una exageración retórica, y merece decirse con precisión porque es la razón
de ser del segundo cálculo.

\begin{matematica}[La $\omega$-regla no sube un escalón: se sale de la escala]
Es un resultado clásico que $\mathrm{PA}^{\omega}$ —PA cerrada bajo la $\omega$-regla— demuestra
\emph{exactamente} las sentencias verdaderas en $\mathbb{N}$:
\[ \mathrm{PA}^{\omega} \;=\; \mathrm{Th}(\mathbb{N}). \]
Y la demostración no usa casi nada de PA: le basta con que los axiomas sean verdaderos en
$\mathbb{N}$ y que las sentencias atómicas cerradas estén decididas. \textbf{Vale igual partiendo
de la aritmética de Robinson}, que es de donde parte este proyecto.
\end{matematica}

\selee{«añadir la $\omega$-regla a una teoría aritmética verdadera la convierte en la teoría de
\emph{todo} lo que es cierto de los números»}

O sea: $\vdash$ no está «un poco por encima de PA». Está en el tope, y por eso no es
\defterm{r.e.}: no hay nada que enumerar cuando ya se demuestra todo lo verdadero. El precio de esa
potencia es que sus demostraciones no son objetos finitos, y por tanto \textbf{no se pueden
codificar} — que es justo lo que la incompletitud necesita hacer.

\begin{muro}[La salvedad que hay que declarar]
El enunciado de arriba es sobre el modelo estándar. El vocabulario de este proyecto no es sólo el
de la aritmética: incluye los 107 axiomas de la capa de codificación, con símbolos opacos cuyo
significado fija la propia teoría. La igualdad $\mathrm{Th}(\mathbb{N})$ se lee, entonces, sobre
$\mathbb{N}$ con esos símbolos interpretados como se pretende. No cambia la moraleja —$\vdash$ no
es r.e., y por eso hace falta el segundo cálculo— pero sí la letra pequeña.
\end{muro}

\section{Cuatro lenguajes objeto, no uno}

Del otro lado pasa algo parecido, y es más fácil de pasar por alto porque los cuatro se llaman
igual. Un \defterm{lenguaje objeto} es sólo una gramática: dice cómo se escriben términos y
fórmulas, y no afirma nada. La teoría aparece cuando se le añaden axiomas.

\begin{center}\small
\begin{tabularx}{\linewidth}{@{}l >{\raggedright\arraybackslash}X >{\raggedright\arraybackslash}X r@{}}
\toprule
estrato & qué añade & para qué sirve & axiomas \\
\midrule
FOL$^=$ desnudo & la gramática y la lógica & escribir fórmulas; no afirma nada & --- \\
Q\texttt{++} & la \defterm{signatura} aritmética y sus \defterm{axiomas objeto} & \textbf{ser incompleta}: es la teoría de la que habla el teorema & 34 \\
Q\texttt{++} codificante & los símbolos opacos del verificador y sus ecuaciones & \textbf{hablar de sintaxis}: sin esto no hay predicado de demostrabilidad & +107 \\
con inducción (\modulo{Full/}) & el esquema de inducción & demostrar D2 y D3 & +1 esquema \\
\bottomrule
\end{tabularx}
\end{center}

\subsection*{Los dos del medio no se parecen en nada, y la diferencia decide el teorema}

Los estratos primero y último se distinguen solos. Los dos de en medio se confunden —los dos se
llaman «Q\texttt{++}», los dos son listas de fórmulas postuladas— y son de naturaleza distinta.

Los \textbf{34} son aritmética: hablan de sumar, multiplicar, comparar. Se entienden mirándolos, se
discuten como se discute un axioma, y son los que hacen del sistema \emph{una teoría de números},
que es lo que el teorema de Gödel exige para tener algo que decir. El
capítulo~\ref{cap:qpp} los desmenuza uno a uno.

Los \textbf{107} no hablan de números: hablan de \emph{sintaxis codificada como} números. Son las
ecuaciones que fijan el comportamiento de símbolos que no se pueden definir —\ident{substfc},
\ident{lenc}, \ident{nthc}, \ident{carc}, el verificador entero— y su papel no es ser verdad sobre
$\mathbb{N}$, sino \textbf{dar a la teoría un vocabulario para hablar de sus propias fórmulas}. Sin
ellos, el predicado «$x$ es demostrable» no se puede ni escribir dentro del lenguaje. El
capítulo~\ref{cap:verificador} los construye.

\begin{observacion}
La proporción es el dato: \textbf{tres veces más maquinaria que aritmética}. Eso es, medido, lo que
cuesta que una teoría pueda hablar de su propia sintaxis — y es la cifra que más sorprende a quien
llega desde la exposición clásica, donde la codificación se despacha con un «y esto se aritmetiza
de manera rutinaria».
\end{observacion}

\begin{muro}[Y por qué la frontera entre los dos no es cosmética]
El verificador consulta la lista de axiomas de la teoría: la regla que acepta «esto es un axioma»
comprueba que el código de la fórmula está en esa lista. Luego \textbf{la lista entra en la
definición del predicado de demostrabilidad}, y ese predicado entra en la sentencia de Gödel. Mover
un axioma de fuera a dentro no hace la teoría «un poco más fuerte»: cambia el predicado, y con él
cambia $G$. Es literalmente otro teorema.

Por eso este libro dice siempre de qué estrato habla, y por eso \adr{015} decidió \emph{definir}
la buena formación de los códigos en vocabulario ya existente en vez de postularla. La historia
completa está en \capfuturo{definir en vez de axiomatizar}.
\end{muro}

\begin{leccion}[Cuatro teorías, cuatro teoremas de incompletitud]
No son cuatro presentaciones de lo mismo: son cuatro teorías distintas, y el teorema de
incompletitud que se demuestra sobre cada una es un teorema distinto. En cuanto la teoría contiene
un verificador que consulta su propia lista de axiomas, \textbf{cambiar la lista cambia la sentencia
de Gödel}. El libro vuelve sobre esto dos veces, y las dos con consecuencias prácticas.
\end{leccion}

\section{Los signos que hacen falta desde ya}

Cuatro notaciones aparecen a partir de aquí constantemente, y conviene fijarlas antes.

\begin{center}\small
\begin{tabularx}{\linewidth}{@{}l >{\raggedright\arraybackslash}X@{}}
\toprule
$\Gamma \vdash A$ & $A$ se deriva de $\Gamma$ en el cálculo $\omega$ \\
$\defnot{\Prf}\,A$ & $A$ se deriva en el cálculo de Hilbert finitario, sin hipótesis
  \selee{«$A$ tiene una demostración finitaria»} \\
$\defnot{\codigo{\varphi}}$ & el \defterm{código de Gödel} de $\varphi$, \textbf{calculado desde
  fuera}: se lo damos ya hecho a la teoría \selee{«el código de $\varphi$»} \\
$\defnot{\dotado{a}}$ & «$a$ dotado»: el código de $a$, \textbf{calculado dentro} por una función
  del propio lenguaje \selee{«$a$ dotado», o «el código de lo que valga $a$»} \\
$\defnot{\Prov}(\codigo{\varphi})$ & la fórmula \emph{objeto} que dice «$\varphi$ es demostrable»
  \selee{«es demostrable», dicho dentro de la teoría} \\
\bottomrule
\end{tabularx}
\end{center}

\noindent Y dos convenios del lenguaje que sorprenden la primera vez. Las variables se escriben
como \defterm{índice de De Bruijn} —\texttt{\#0} es la ligada más cercana, y los cuantificadores no
llevan nombre de variable—. Y el término que denota el número $n$ es el \defterm{numeral}
$\sigma^n 0$: es \emph{canónico}, sólo hay una forma de escribirlo, y de ahí que sirva para
representar códigos.

\subsection*{Dos maneras de decir «el código de», y no son intercambiables}

Las dos filas segunda y tercera parecen decir lo mismo y no lo dicen. La diferencia es \emph{quién
hace el cálculo}, y de ella depende medio libro.

$\codigo{\varphi}$ es una \textbf{función de Lean}: nosotros, desde fuera, tomamos una fórmula
concreta y calculamos su código. El resultado es un término cerrado y concreto que se puede escribir
en un papel. Para usarlo hay que tener la fórmula \textbf{en la mano}.

$\dotado{a}$ es un \textbf{símbolo de función del lenguaje objeto} —se llama \ident{tcFn}— que la
teoría aplica a lo que sea. No calcula nada por adelantado: es una expresión del lenguaje, como
$\obj{\sigma}$ o $\obj{+}$.

\begin{center}\small
\begin{tabularx}{\linewidth}{@{}l >{\raggedright\arraybackslash}X >{\raggedright\arraybackslash}X@{}}
\toprule
 & $\codigo{\cdot}$ (\ident{formCode}, \ident{termCode}) & $\dotado{\cdot}$ (\ident{tcFn}) \\
\midrule
quién lo calcula & nosotros, en la metateoría & la teoría, dentro del lenguaje \\
qué admite como argumento & una fórmula o término \textbf{concretos} & cualquier término, incluida \textbf{una variable} \\
qué es el resultado & un término cerrado, escrito & una expresión con variables libres \\
¿se puede poner bajo un $\forall$? & \textbf{no} útilmente & \textbf{sí} \\
\bottomrule
\end{tabularx}
\end{center}

\noindent La cuarta fila es la que importa, y se ve en un ejemplo. Ésta es una de las ecuaciones que
caracterizan a \ident{tcFn}, y dice, para \emph{todo} $x$, cómo se relaciona el código de
$\obj{\sigma}x$ con el de $x$:
\[ \forall x.\quad \dotado{(\obj{\sigma}x)} \;\doteq\; \langle 1,\ \codigo{\obj{\sigma}},\ [\,\dotado{x}\,]\rangle \]
\selee{para todo $x$, el código del sucesor de $x$ es la lista formada por la etiqueta de
aplicación, el código del símbolo $\sigma$, y el código de $x$.}

\begin{muro}[Por qué ahí no se puede escribir $\codigo{x}$]
Se podría \emph{teclear}, y ahí está la trampa: $\codigo{x}$ es legal, pero significa «el código del
término \emph{la variable~$x$}» — un código concreto, cerrado, de un objeto sintáctico. \textbf{No}
significa «el código de lo que valga $x$», que es lo que la fórmula necesita. El punto es la única
notación que puede ir bajo un cuantificador diciendo lo segundo.
\end{muro}

\begin{leccion}[Y una asimetría que se cobra cara]
Como $\dotado{\cdot}$ es un símbolo de función, la teoría demuestra que \textbf{respeta la
igualdad}: de $a \doteq b$ se sigue $\dotado{a} \doteq \dotado{b}$. $\codigo{\cdot}$ no tiene nada
parecido, porque no es un símbolo del lenguaje: es una operación sobre \emph{sintaxis}.

Y ahí está el precio. Un símbolo de función sólo puede depender del \textbf{valor} de su argumento,
nunca de cómo se escribió ese valor. Pero un mismo número es a la vez un numeral y una lista —
$\obj{\mathrm{cons}}(\mathrm{nil},\mathrm{nil})$ y $\obj{\sigma\sigma}\obj{0}$ son el mismo número—,
así que exigirle a $\dotado{\cdot}$ que devuelva un código distinto según «cómo se escribió» es
pedirle lo imposible. Eso hizo la teoría inconsistente durante meses, y es el episodio central de la
Parte~\ref{parte:noescrito}.
\end{leccion}

\section{Cómo se lee la chapa}

De todo lo anterior sale el instrumento que este libro usa en cada fragmento de código y en cada
enunciado: una \textbf{chapa de dos casillas} que dice quién afirma y sobre qué. No es un adorno —
es la distinción cuya violación costó 31 módulos, y \textbf{no es binaria}: son dos ejes
independientes.

%% La leyenda de niveles. Va ANTES del índice: PLAN-LIBRO.md §2.5.
\section*{2 · Los dos ejes: quién afirma, y sobre qué}

Con esas palabras ya se puede plantear la pregunta que vuelve en cada línea:
\emph{¿esto es lenguaje objeto o lenguaje meta?} En un libro de metamatemática se
plantea en cada línea, y responderla mal es caro: la inconsistencia que cuenta la
Parte~IV fue exactamente un error de categoría. Por eso aquí el nivel \textbf{se imprime}.

No es una distinción binaria. Hay dos ejes independientes, y cada fragmento de código y cada
enunciado llevan una chapa con los dos:

\noindent\textbf{Eje 1 — ¿quién afirma?} Son cuatro relaciones de derivabilidad distintas, y
\emph{de tipos distintos}: lo que las separa es si arrastran contexto y si son recursivamente
enumerables.

\begin{center}\footnotesize\setlength{\tabcolsep}{4pt}
\begin{tabularx}{\linewidth}{@{}l >{\ttfamily\scriptsize}l cc >{\raggedright\arraybackslash}X@{}}
\toprule
chapa & tipo en Lean & ctx. & r.e. & papel \\
\midrule
\colorbox{nvLean}{\textcolor{white}{\sffamily\scriptsize\bfseries\strut\,Lean\,}} &
  \normalfont--- & --- & --- & la metateoría: un teorema \emph{sobre} el sistema \\
\colorbox{nvDer}{\textcolor{white}{\sffamily\scriptsize\bfseries\strut\,$\Gamma\vdash$\,}} &
  List Formula → Formula → Prop & sí & \textbf{no} & el cálculo $\omega$; sólido en $\mathbb{N}$ \\
\colorbox{nvPrf0}{\textcolor{white}{\sffamily\scriptsize\bfseries\strut\,Prf$_0$\,}} &
  Formula → Prop & no & sí & Hilbert \textbf{intuicionista} \\
\colorbox{nvPrf}{\textcolor{white}{\sffamily\scriptsize\bfseries\strut\,Prf\,}} &
  Formula → Prop & no & sí & Hilbert clásico: \textbf{el que se aritmetiza} \\
\colorbox{nvPrfH}{\textcolor{white}{\sffamily\scriptsize\bfseries\strut\,$\Gamma\,$PrfH\,}} &
  List Formula → Formula → Prop & sí & sí & variante con contexto; ahí vive la deducción \\
\bottomrule
\end{tabularx}
\end{center}

\noindent\textbf{Eje 2 — ¿sobre qué?} Cinco valores.

\begin{center}\small
\begin{tabularx}{\linewidth}{@{}l >{\raggedright\arraybackslash}X@{}}
\toprule
\colorbox{nvLean!18}{\textcolor{nvLean}{\sffamily\scriptsize\bfseries\strut\,Lean\,}} &
  objetos nativos: \ident{Nat}, listas, recursión estructural \\
\colorbox{nvLean!18}{\textcolor{nvLean}{\sffamily\scriptsize\bfseries\strut\,sintaxis\,}} &
  la representación en Lean de la sintaxis: \ident{Term}, \ident{Formula} \\
\colorbox{nvObj!18}{\textcolor{nvObj}{\sffamily\scriptsize\bfseries\strut\,objeto\,}} &
  enunciados de la teoría aritmética: $\obj{+}$, $\obj{\cdot}$, $\obj{\sigma}$, sus axiomas \\
\colorbox{nvCod!18}{\textcolor{nvCod}{\sffamily\scriptsize\bfseries\strut\,código\,}} &
  términos objeto que \emph{denotan} sintaxis: $\cod{\codigo{\varphi}}$, \ident{numeral} \\
\colorbox{nvProv!18}{\textcolor{nvProv}{\sffamily\scriptsize\bfseries\strut\,Prov\,}} &
  dentro del predicado de demostrabilidad: la teoría hablando de sí misma \\
\bottomrule
\end{tabularx}
\end{center}

\noindent \textbf{Los puentes, y el que no existe.} Las cuatro no son intercambiables, y la
dirección importa:
\[
  \Prf_0\,\varphi \;\longrightarrow\; \Prf\,\varphi \;\longrightarrow\;
  \bigl(\forall \Gamma,\ \mathsf{PrfH}\ \Gamma\ \varphi\bigr),
  \qquad
  \Prf\,\varphi \;\longrightarrow\; \text{\ident{axioms}} \vdash \varphi .
\]
El recíproco $\text{\ident{axioms}} \vdash \varphi \to \Prf\,\varphi$ \textbf{es falso, y tiene que
serlo}: si valiera, $\vdash$ sería r.e.\ y Tarski cerraría el paso. Ésa es la asimetría que obliga a
tener dos cálculos en vez de uno, y por eso la chapa distingue en cuál se afirma.
Además \ident{prf0_to_derives} no usa \emph{ninguna} meta-regla $\omega$ y \ident{prf_to_derives}
usa \ident{dne} exactamente una vez: la diferencia entre los dos puentes es, medida,
$\{\ident{FOL.MetaRules.dne}\}$.

\medskip\noindent Tres ejemplos que ocupan tres casillas distintas y que conviene comparar:
\ident{two_mul_consN} afirma \emph{en Lean} sobre \ident{Nat};
\ident{prf_formCode_numeral} afirma \emph{en Prf} sobre \textbf{códigos};
\ident{hFN} afirma lo mismo, pero \emph{en} $\vdash$.
Y \ident{pcc_dot_cons} afirma en \ident{Prf} \textbf{dentro de} \ident{Prov} — cuatro niveles de
distancia entre el primero y el último.

\begin{leccion}[Por qué la chapa es un control, no un adorno]
La cuarta fila del segundo eje es la que se cobró 31 módulos. \ident{tcFn} —«el código del
código»— es una operación sobre \textbf{sintaxis} que se declaró como función \textbf{objeto}, y
una función objeto sólo puede depender de \textbf{valores}. Con la chapa puesta, la mezcla se ve
antes de escribirla.
\end{leccion}


\begin{leccion}[La distinción que hay que llevarse de aquí]
$\Prov$ es una fórmula del lenguaje objeto; \defterm{provable} quiere decir que la propia teoría
demuestra algo, no que lo comprobemos nosotros desde fuera. La diferencia entre comprobar una cosa y
que la teoría la demuestre \emph{de sí misma} es toda la distancia que hay entre la primera condición
de derivabilidad y la tercera — es decir, entre lo que está hecho y lo que falta.
\end{leccion}
